\documentclass[../../main.tex]{subfiles}
\begin{document}

{\bf [解答]}

与えられた数列
\begin{align}
  \begin{cases}
    a_1 \le a_2 \le \cdots \le a_n \\
    a_1 + a_2 + \cdots + a_n = 0
  \end{cases} \label{eq:1}
\end{align}
について，第一式より
\begin{align}
  & n a_1 \le a_1 + a_2 + \cdots + a_n \le n a_n \\
  \therefore
  & a_1 \le 0 \le a_n \quad (\because n \in \mathbb{N}, \text{\cref{eq:1}})
\end{align}
となる．
従って$a_k \le 0 \le a_{k+1} \ (k=1, 2, \cdots, n-1)$ をみたす$k$が存在する．
この時，示すべき不等式の左辺は
\begin{align}
  \sum_{i=1}^{n} i a_i
  &= (a_1 + \cdots + a_n) + (a_2 + \cdots + a_n) + \cdots + (a_n) \\
  &= 0 + (-a_1) + (-a_1 - a_2) + \cdots + (-a_1 - \cdots - a_k) \\
  &\quad + (a_{k+1} + \cdots + a_n) + \cdots + a_n \\
  &> 0 \quad (\because \text{各項0以上，等号成立は} a_1 = a_2 = \cdots = a_n = 0 \text{だが，題意よりそれはない})
\end{align}
となって題意は示された．

{\bf [別解]}
数学的帰納法でも示せる．簡単に証明の概要を示す．
$n=k$での成立を仮定して$a_1, a_2, \cdots, (a_k + a_{k+1})$に適用すると
\begin{align}
  0 &< \sum_{i=1}^{k} i a_i + k a_{k+1} \\
  &\le \sum_{i=1}^{k+1} i a_i \quad (\because a_{k+1} \ge 0)
\end{align}
となって示される．

\end{document}
